\documentclass[11pt]{article}
\usepackage[margin=1in]{geometry}
\usepackage{enumitem}
% =============================================================================
% Preambles/header.tex — shared LaTeX/Beamer preamble for this project
%
% Use from a Beamer lecture by prepending:
%     \input{header}
% and compiling with TEXINPUTS=../Preambles:$TEXINPUTS (the /compile-latex
% skill does this automatically).
%
% PALETTE CONTRACT. Color names defined here MUST match the names in
% Quarto/theme-template.scss so Beamer and Quarto renderings use the same
% palette. The scripts/check-palette-sync.sh script enforces this.
%
% To customize: edit the HEX values below AND the matching $variable in
% Quarto/theme-template.scss, then run scripts/check-palette-sync.sh.
% =============================================================================

% -----------------------------------------------------------------------------
% Engine detection + encoding
%
% Our /compile-latex skill uses XeLaTeX, where inputenc is unnecessary and
% can produce encoding warnings. Only load inputenc under pdfTeX.
% -----------------------------------------------------------------------------
\usepackage{iftex}
\ifPDFTeX
  \usepackage[utf8]{inputenc}
\fi

\usepackage{amsmath, amssymb, amsthm}
\usepackage{booktabs}
\usepackage{graphicx}

% -----------------------------------------------------------------------------
% PALETTE — mirrors Quarto/theme-template.scss variable names.
% Load xcolor and define colors BEFORE anything that references them
% (notably hyperref, hypersetup, and the Beamer color assignments).
% -----------------------------------------------------------------------------
\usepackage{xcolor}

\definecolor{primary-blue}{HTML}{012169}      % headings, accents
\definecolor{primary-gold}{HTML}{B9975B}      % emphasis, borders
\definecolor{highlight-yellow}{HTML}{F2A900}  % markers, alerts
\definecolor{light-bg}{HTML}{E8EDF5}          % title / section backgrounds
\definecolor{jet}{HTML}{1A1A1A}               % body text

% Semantic colors (used in diagrams and annotations)
\definecolor{positive}{HTML}{15803D}          % good / observed / identified
\definecolor{negative}{HTML}{B91C1C}          % bad / problematic / bias
\definecolor{neutral}{HTML}{525252}           % reference / context

% Highlight accents (match .hi-* classes in the SCSS)
\definecolor{hi-slate}{HTML}{314F4F}
\definecolor{hi-green}{HTML}{2E8B57}
\definecolor{hi-red}{HTML}{C41E3A}

% Semantic aliases so skill templates and snippets can write
% \textcolor{accent}{...} without hard-coding "primary-blue".
\colorlet{accent}{primary-blue}
\colorlet{accent2}{primary-gold}

% -----------------------------------------------------------------------------
% Hyperref — loaded AFTER xcolor + palette so link colors resolve.
% -----------------------------------------------------------------------------
\usepackage{hyperref}
\hypersetup{colorlinks=true, linkcolor=primary-blue, urlcolor=primary-blue,
            citecolor=primary-blue, pdfborder={0 0 0}}

% -----------------------------------------------------------------------------
% Course code — set once per deck (after \input{header}, before \begin{document})
% via \coursecode{CS401}. Shown in the footer on every slide so a deck's course
% is identifiable without opening the title page. Empty by default.
% -----------------------------------------------------------------------------
\makeatletter
\newcommand{\coursecode}[1]{\gdef\@coursecodetext{#1}}
\gdef\@coursecodetext{}
\makeatother

% -----------------------------------------------------------------------------
% Beamer theme assignments (only apply under Beamer).
%
% We detect Beamer via \@ifclassloaded{beamer}, which is the documented,
% non-internal way. This must be wrapped in \makeatletter/\makeatother
% because \@ifclassloaded contains an @-catcode character.
% -----------------------------------------------------------------------------
\makeatletter
\@ifclassloaded{beamer}{%
  \usefonttheme{professionalfonts}
  \setbeamercolor{structure}{fg=primary-blue}
  \setbeamercolor{frametitle}{fg=primary-blue}
  \setbeamercolor{title}{fg=primary-blue}
  \setbeamercolor{subtitle}{fg=primary-gold}
  \setbeamercolor{author}{fg=primary-blue}
  \setbeamercolor{institute}{fg=neutral}
  \setbeamercolor{date}{fg=primary-gold}
  \setbeamercolor{itemize item}{fg=primary-blue}
  \setbeamercolor{itemize subitem}{fg=primary-gold}
  \setbeamercolor{alerted text}{fg=primary-gold}
  \setbeamercolor{block title}{fg=white, bg=primary-blue}
  \setbeamercolor{block body}{bg=light-bg}
  % Semantic block variants: block=definition/concept (blue), exampleblock=
  % worked example (green/positive), alertblock=key takeaway or caution (gold).
  % Keep to <=2 colored boxes per slide across ALL three variants (INV-7).
  \setbeamercolor{block title example}{fg=white, bg=positive}
  \setbeamercolor{block body example}{bg=light-bg}
  \setbeamercolor{block title alerted}{fg=white, bg=primary-gold}
  \setbeamercolor{block body alerted}{bg=light-bg}

  % Minimal footer: course code (left) + page X / N (right), no navigation clutter
  \setbeamertemplate{navigation symbols}{}
  \setbeamertemplate{footline}{%
    \color{neutral}\footnotesize\hspace{0.7em}\@coursecodetext\hfill
    \insertframenumber/\inserttotalframenumber\hspace{0.7em}\vspace{0.3em}}
}{}
\makeatother

% -----------------------------------------------------------------------------
% TikZ setup — libraries the snippet gallery uses
% -----------------------------------------------------------------------------
\usepackage{tikz}
\usetikzlibrary{arrows.meta, positioning, calc, decorations.pathreplacing,
                fit, shapes.geometric, backgrounds}

% Shared TikZ styles — snippets and hand-written diagrams can pick these up
% via `\begin{tikzpicture}[every node/.style={...}]` or by naming specific
% styles (e.g., `\node[dag-node] (X) at (0,0) {X};`).
\tikzset{
  % Node styles for causal DAGs and similar
  dag-node/.style = {
    draw=primary-blue, fill=light-bg, rounded corners=2pt,
    minimum width=1.2cm, minimum height=0.9cm, align=center, font=\small
  },
  decision-node/.style = {
    draw=primary-gold, fill=white, diamond, aspect=1.8,
    minimum width=1.8cm, minimum height=1.2cm, align=center, font=\small,
    inner sep=1pt
  },
  % Edge styles — observed vs counterfactual
  observed-edge/.style = {-{Latex[length=2.5mm]}, thick, draw=jet},
  counterfactual-edge/.style = {-{Latex[length=2.5mm]}, thick, draw=neutral, dashed},
  confound-edge/.style = {-{Latex[length=2.5mm]}, thick, draw=negative, dashed},
  % Data-point styles for plots
  observed-dot/.style = {fill=primary-blue, draw=primary-blue, circle, inner sep=1.2pt},
  counterfactual-dot/.style = {fill=white, draw=neutral, circle, inner sep=1.2pt}
}

% -----------------------------------------------------------------------------
% Convenience macros
% -----------------------------------------------------------------------------
\newcommand{\muted}[1]{\textcolor{neutral}{#1}}
\newcommand{\key}[1]{\textcolor{primary-gold}{\textbf{#1}}}
\newcommand{\good}[1]{\textcolor{positive}{#1}}
\newcommand{\bad}[1]{\textcolor{negative}{#1}}

% Standout frame macro: full-bleed dark background for section transitions.
% Beamer-only (uses \begin{frame}/\setbeamercolor/\usebeamerfont) -- do not
% call from an article-class file (e.g. Notes/<CODE>/*-notes.tex). \muted,
% \key, \good, \bad, and \coursecode above are plain xcolor/gdef and are
% safe to use from either class; verified when Notes/article-class support
% was added.
% Expand: \transitionslide{Your transition title}
\newcommand{\transitionslide}[1]{%
  \begingroup
  \setbeamercolor{background canvas}{bg=primary-blue}
  \begin{frame}[plain]
    \centering
    \vfill
    {\usebeamerfont{title}\color{white}\Large #1\par}
    \vfill
  \end{frame}
  \endgroup
}

% Section divider frame (Beamer-only), an alternative to \transitionslide with a
% darker two-line style: near-black (jet) background, a small white label line
% above a larger gold title, and a thin gold rule along the bottom edge.
% Expand: \sectiondivider{Section 2}{Pointers}
\newcommand{\sectiondivider}[2]{%
  \begingroup
  \setbeamercolor{background canvas}{bg=jet}
  \begin{frame}[plain]
    \vspace*{3.0cm}
    {\color{white}\ttfamily\large #1\par}
    \vspace{0.5cm}
    {\color{primary-gold}\LARGE #2\par}
    \begin{tikzpicture}[remember picture, overlay]
      \draw[primary-gold, line width=2.5pt]
        ($(current page.south west)+(0,0.1)$) -- ($(current page.south east)+(0,0.1)$);
    \end{tikzpicture}
  \end{frame}
  \endgroup
}

% -----------------------------------------------------------------------------
% Channel watermark — "Concepts That Click" (YouTube).
%
% Article-class files (CompetitiveExam Q/A sets, Notes, Assignments) get a
% light diagonal draft watermark on every page plus a centered footer naming
% the channel. Beamer decks get a subtle TikZ background watermark instead
% (the footline already carries the course code). Centralized here so every
% MCQ PDF is branded without per-file edits.
% -----------------------------------------------------------------------------
\makeatletter
\@ifclassloaded{article}{%
  \usepackage{draftwatermark}
  \SetWatermarkText{Concepts That Click}
  \SetWatermarkScale{1}
  \SetWatermarkFontSize{2cm}
  \SetWatermarkLightness{0.86}
  \SetWatermarkAngle{45}
  \usepackage{fancyhdr}
  \pagestyle{fancy}
  \fancyhf{}
  \fancyfoot[C]{\footnotesize\textcolor{neutral}{Concepts That Click~|~YouTube: \href{https://www.youtube.com/@concepts.thatclick}{@concepts.thatclick}}}
  \renewcommand{\headrulewidth}{0pt}
  \renewcommand{\footrulewidth}{0pt}
  \fancypagestyle{plain}{%
    \fancyhf{}%
    \fancyfoot[C]{\footnotesize\textcolor{neutral}{Concepts That Click~|~YouTube: \href{https://www.youtube.com/@concepts.thatclick}{@concepts.thatclick}}}%
    \renewcommand{\headrulewidth}{0pt}%
    \renewcommand{\footrulewidth}{0pt}%
  }
}{}
\@ifclassloaded{beamer}{%
  \setbeamertemplate{background}{%
    \begin{tikzpicture}[remember picture, overlay]
      \node[rotate=45, opacity=0.055, align=center]
        at (current page.center)
        {\Huge\textcolor{neutral}{Concepts That Click}};
    \end{tikzpicture}%
  }
}{}
\makeatother 
\coursecode{JKSSB}
\renewcommand{\thesection}{39.\arabic{section}}
\newtheorem{example}{Example}[section]
\newenvironment{definitionbox}[1]{%
  \par\noindent\textbf{Definition (#1).}\ \itshape
}{\par\vspace{0.3em}}
\title{E-Banking --- Detailed Lecture Notes}
\author{JKSSB Computer Awareness}
\date{\today}

\begin{document}
\maketitle

\section{What e-banking covers}
Electronic banking, or e-banking, means using banking services without
visiting the branch. The account itself stays in the bank; only the channel
used to reach it changes. \textbf{Net banking} is banking through the bank's
website in a browser. \textbf{Mobile banking} is banking through the bank's
app on the phone. \textbf{UPI} (Unified Payments Interface) apps move money
directly between bank accounts from the phone. An ATM, the net-banking site,
the mobile app, and a UPI app are therefore four doors into the same
account.

\begin{definitionbox}{E-banking}
The use of banking services through electronic channels --- net banking,
mobile banking, UPI, and ATMs --- instead of visiting the branch in person.
\end{definitionbox}

\begin{center}
\begin{tabular}{p{0.24\textwidth}p{0.66\textwidth}}
\toprule
\textbf{Term} & \textbf{Meaning in this lesson} \\
\midrule
Net banking & Banking through the bank's website in a browser. \\
Mobile banking & Banking through the bank's official app. \\
UPI & Unified Payments Interface; instant phone-based payments between accounts. \\
VPA & Virtual Payment Address (for example name@bank); identifies the account without the account number. \\
ATM PIN & Personal Identification Number used with the physical card. \\
UPI PIN & Personal Identification Number authorising UPI payments. \\
CVV & Card Verification Value printed on the card. \\
OTP & One-Time Password sent by SMS, valid for one transaction only. \\
MFA & Multi-Factor Authentication; checking two or more factors together. \\
\bottomrule
\end{tabular}
\end{center}

The rest of this lesson uses these terms in one consistent sense. A
\textbf{PIN} is a secret typed by the holder. \textbf{CVV} is printed on the
card. \textbf{OTP} arrives by SMS and expires. \textbf{MFA} is the
combination of two or more different proofs of identity.

\section{UPI in detail}
UPI stands for Unified Payments Interface. Each user is identified by a
\textbf{VPA} (Virtual Payment Address), such as name@bank, so money can be
sent without sharing the account number or IFSC. Transfers are
\textbf{instant} and work \textbf{24 $\times$ 7}, including Sundays and
holidays. Every payment is authorised on the payer's own phone by entering
the \textbf{UPI PIN}, which is set separately for each bank account inside
the UPI app. The central rule of UPI safety is short: entering the UPI PIN
always sends money, and receiving money never needs any PIN.

Other electronic transfer modes appear in nearby options and should be kept
in one line each. \textbf{IMPS} (Immediate Payment Service) is an instant
bank-to-bank transfer available 24 $\times$ 7. \textbf{NEFT} (National
Electronic Funds Transfer) settles in batches and has no minimum amount.
\textbf{RTGS} (Real Time Gross Settlement) settles large-value transfers
instantly and carries a minimum amount. UPI and IMPS are the instant,
phone-friendly modes; NEFT is batched; RTGS is for large values.

\begin{definitionbox}{UPI}
Unified Payments Interface: an instant, round-the-clock system for moving
money between bank accounts from a phone, identified by VPA and authorised by
a UPI PIN.
\end{definitionbox}

\section{The four secret codes}
Four codes look alike and are deliberately mixed in distractors. Each lives
in a different place and does a different job.

\textbf{ATM PIN.} The Personal Identification Number used with the physical
card at an ATM or shop terminal. It is memorised, typed on the machine, and
authorises card transactions on that channel. It never authorises a UPI
payment.

\textbf{UPI PIN.} The Personal Identification Number typed in the phone to
authorise a UPI payment. It is set per bank account in the UPI app and can
differ from the ATM PIN. Approving any screen with the UPI PIN sends money.

\textbf{CVV.} The Card Verification Value: the 3-digit number printed on the
back of the card (4 digits on the front for some cards). It is entered with
the card number and expiry date for online card payments and proves the
card is in hand. It is fixed and reusable for the life of the card.

\textbf{OTP.} The One-Time Password: a short numeric code sent by SMS that
is valid for one transaction only and for a few minutes. It is generated
fresh each time and expires after use or after the time window.

\begin{center}
\begin{tabular}{p{0.16\textwidth}p{0.24\textwidth}p{0.22\textwidth}p{0.28\textwidth}}
\toprule
\textbf{Code} & \textbf{Where it lives} & \textbf{How long it lasts} & \textbf{Typical use} \\
\midrule
ATM PIN & In memory & Until changed & Card at ATM or shop terminal \\
UPI PIN & In memory & Until changed & Authorising a UPI payment \\
CVV & Printed on card & Life of the card & Online card payment with card number \\
OTP & Sent by SMS & One transaction, minutes & Second approval for a login or payment \\
\bottomrule
\end{tabular}
\end{center}

\begin{example}
A question states, ``The 3-digit number on the back of the card was sent by
SMS and expired after one use.'' This mixes two codes. The number printed on
the back of the card is the \textbf{CVV}, which is fixed and reusable; the
code that arrives by SMS and expires after one use is the \textbf{OTP}. The
stem has attached OTP behaviour to the CVV name.
\end{example}

\section{Authentication and multi-factor authentication}
\textbf{Authentication} means proving ``I am the account holder'' before the
bank acts. It uses three kinds of factor. \textbf{Something you know} is a
password or PIN. \textbf{Something you have} is the phone, the card, or the
OTP that arrives on the registered mobile number. \textbf{Something you are}
is a fingerprint or face scan.

\begin{definitionbox}{Multi-Factor Authentication (MFA)}
Authentication that checks two or more different factors together, for
example a password (know) plus an SMS OTP (have), or a PIN plus a
fingerprint.
\end{definitionbox}

A login password plus an SMS OTP is therefore \textbf{two-factor}
authentication, not one: the password proves knowledge and the OTP proves
possession of the registered phone. Asking for the OTP a second time does
not add a new factor; a fingerprint or card would. The exam trap is to call a
password-plus-OTP login ``single-factor'' or to call two passwords
``multi-factor''; factors must be of \emph{different} kinds.

\section{QR codes and collect requests}
A \textbf{QR} (Quick Response) code is the square barcode the merchant
displays. Scanning it with the UPI app fills in the merchant's payment
address; the payer then checks the \textbf{merchant name and amount} on
their own screen and approves with the UPI PIN. A QR never asks for a PIN or
OTP by itself. A common fraud is a fake QR sticker pasted over the real one,
so the displayed name on the payer's screen is the check that matters: if
the name does not match the shop, stop.

A \textbf{collect request} is a message asking the receiver to pay: it
appears as ``pay Rs.\ X to Y, approve with UPI PIN''. Approving it and
entering the UPI PIN sends the receiver's money to the requester. No refund,
cashback, prize, or buyer ever needs the seller to enter a UPI PIN; entering
the PIN always means paying. Unexpected collect requests, especially ones
with remarks like ``refund'' or ``cashback'', must be declined or ignored.
To receive money, the receiver only shares a VPA or QR and enters nothing.

\begin{example}
A buyer sends a Rs.\ 500 collect request to a seller with the remark
``advance, approve to receive''. The seller's screen shows ``pay Rs.\ 500'',
and approving with the UPI PIN would send Rs.\ 500 to the buyer. The remark
is false: receiving needs no PIN. The safe action is to decline and ask the
buyer to send the payment instead.
\end{example}

\section{Safe-banking rules}
The rules below are tested almost verbatim, so they are stated exactly.

\begin{enumerate}
  \item The bank \textbf{never} asks for OTP, CVV, ATM PIN, or UPI PIN by
    call, SMS, email, or link. Any such request is fraud.
  \item \textbf{Never share} an OTP for any reason --- to ``cancel'', to
    ``verify'', or to ``receive'' a payment. Sharing it completes the
    transaction the fraudster started.
  \item Do not open \textbf{KYC-update or account-blocked} links from SMS.
    Open the bank's site or app by typing the address or using the official
    app.
  \item \textbf{Never install} a screen-sharing or remote-access app at a
    caller's request; it exposes OTPs, PINs, and balances on the screen.
  \item Use \textbf{typed addresses and official apps}, and check the
    \textbf{padlock (https)} on the banking site before logging in.
  \item Avoid \textbf{public Wi-Fi} for banking; prefer mobile data or a
    trusted network.
  \item \textbf{Log out} after net banking, lock the UPI app with the device
    lock, and keep \textbf{SMS and email alerts on} so an unknown debit is
    noticed at once.
\end{enumerate}

\section{Exam retrieval and common confusions}
Six distinctions prevent most errors on this topic.
\begin{enumerate}
  \item ATM PIN is for the \textbf{card} channel; UPI PIN is for the
    \textbf{UPI} channel. They are set and used separately.
  \item CVV is \textbf{printed on the card} and reusable; OTP is
    \textbf{sent by SMS} for one transaction and expires.
  \item Entering the UPI PIN \textbf{always pays}; receiving money needs
    \textbf{no PIN}.
  \item A password plus an SMS OTP is \textbf{two-factor} authentication
    (know + have); two passwords are still one factor.
  \item The bank \textbf{never asks} for OTP, CVV, or any PIN. Any caller
    who asks is not the bank.
  \item A QR is verified by the \textbf{name and amount on your own
    screen}; a pasted-over sticker is defeated by that check.
\end{enumerate}

\begin{example}
An item asks what is needed to receive a UPI refund. Trace the rule:
receiving needs a VPA or QR and no PIN; any screen asking for the UPI PIN is
a payment screen. The answer is that no PIN or OTP is entered to receive;
the ``approve to receive'' request is a collect-request fraud.
\end{example}

\subsection*{Exercises}
\begin{enumerate}[label=\arabic*.]
  \item Expand UPI, OTP, CVV, PIN, QR, and MFA.
  \item Distinguish the ATM PIN from the UPI PIN: where is each used?
  \item Distinguish CVV from OTP: where does each live, and how long does
    each last?
  \item Explain why a login password plus an SMS OTP counts as two-factor
    authentication.
  \item A seller receives a collect request marked ``refund, approve to
    receive''. What should the seller do, and why?
\end{enumerate}

\subsection*{Solutions}
\begingroup\small
\begin{enumerate}[label=\arabic*.]
  \item UPI --- Unified Payments Interface; OTP --- One-Time Password; CVV
    --- Card Verification Value; PIN --- Personal Identification Number; QR
    --- Quick Response code; MFA --- Multi-Factor Authentication.
  \item The ATM PIN is typed on a machine with the physical card at an ATM
    or shop terminal. The UPI PIN is typed in the phone to authorise a UPI
    payment, set per bank account in the UPI app. Neither works on the
    other's channel.
  \item CVV is printed on the card and is reusable for the life of the card;
    it is entered with the card number for online card payments. OTP is sent
    by SMS, is valid for one transaction only and a few minutes, and then
    expires.
  \item The password is something the user knows and the OTP on the
    registered phone is something the user has. Two different factor kinds
    are checked, so the login is two-factor (a form of MFA).
  \item Decline or ignore it. A collect request approved with the UPI PIN
    sends the seller's money; receiving a refund needs no PIN. The remark is
    a fraud pattern.
\end{enumerate}
\endgroup
\end{document}
